% 本文件由 LaTeX 排版 Agent 根据有效 translation.json 自动生成。
% author=Justin Martyr; work_id=0131; title=论复活

\chapter{论复活}

% structure: p0001 source=h1 target=omitted reason=page-title-already-rendered-by-parent

% structure: p0002 source=h2 target=section reason=relative-heading-rank; base=h2

\section{第一章 真理自明之德}

真理之言是自由的，且带有自身的权威，不屑于落入任何精巧的论证，也不愿忍受听众的逻辑审查。但它因其自身的高贵以及差遣祂的那位所应得的信心而被相信。如今，真理之言是由天主派遣的；因此真理所主张的自由并非傲慢。因为它带着权威被派遣，就不应被要求为所言之事提供证明；因为没有任何超越它自身的证明，而它本身就是天主。因为每一个证明都比它所证明的更有力、更可信；因为那不被相信的事，在证明提出之前，一旦证明提出，便获得信任，并被承认为当初所陈述的那样。但没有任何事物比真理更有力或更可信；所以要求证明真理的人，就像希望证明为何感官所显现的事物会显现的人一样。因为通过理性所接受的事物，其检验标准是感官；但感官本身却没有超越自身的检验。因此，就像我们将理性所追求的事物带到感官面前，由它来判断这些事物的性质，无论所说的是真是假，然后就不再审判，完全信任它的决定；同样，我们将所有关于人类和世界的话都引向真理，由它来判断是否毫无价值。但真理的宣告，我们不用任何别的检验标准来判断，完全信任它本身。而天主，宇宙的圣父，完美理智，就是真理。圣言，祂的圣子，来到我们中间，穿上了肉身，启示了祂自己和圣父，在祂自己内赐给我们从死者中的复活，以及随后的永生。这就是耶稣基督，我们的救主和主。因此，祂本身就是对自己的信德和对万物的证明。因此，那些跟随祂、认识祂，并将信德寄托于祂作为证明的人，必安息于祂内。但由于仇敌不懈地抵抗许多人，并使用许多不同的诡计来陷害他们，为诱使有信德者离开他们的信德，并阻止无信德者相信，在我看来，我们也必须拿起信德无敌的道理武装自己，为弱者与它作战。\par

% structure: p0004 source=h2 target=section reason=relative-heading-rank; base=h2

\section{第二章 反对肉身复活的异议}

那些持错误意见的人说，没有肉身复活；他们的理由是，腐败和分解的东西不可能恢复原状。而且除了不可能之外，他们还说肉身的救恩是无益的；他们辱骂肉身，指出它的软弱，并宣称只有它是我们犯罪的起因，因此，他们说，如果肉身复活，我们的软弱也随着它复活。他们又提出这样的诡辩：若肉身复活，它必须要么完整并拥有所有部分，要么不完美。但若它不完美地复活，则表明天主方面能力不足，如果有些部分能得救，而其他部分不能；但若所有部分都得救，那么身体显然就有所有肢体。但说这些肢体在死后复活后仍然存在，岂不荒谬？因为救主说过：\BibleQuote{他们不娶，也不嫁，却像天上的天使一样。}\BibleRef{谷 12:25}他们说，天使既没有肉身，也不吃，也不性交；因此没有肉身复活。他们用这些和类似的论点，试图使人偏离信德。还有一些人主张，连耶稣自己也只是以灵性方式显现，并非有肉身，只是现出肉身的外貌罢了：这些人试图剥夺肉身的应许。那么，首先，让我们解决那些在他们看来无法解决的问题；然后我们按顺序引入关于肉身的论证，证明它分享救恩。\par

% structure: p0006 source=h2 target=section reason=relative-heading-rank; base=h2

\section{第三章 若肢体复活，它们是否仍需履行现在的功能？}

他們說，如果身體將完整復活，並擁有所有肢體，那麼必然隨之而來的是肢體的功能也將存在；子宮將懷孕，男性也將履行其生殖功能，其餘肢體也同樣。現在，讓這個論點取決於這一斷言。因為若此斷言被證明為假，他們的整個異議將被消除。現在，很明顯，那些履行功能的肢體所履行的功能是我們在今生所看到的，但這並不意味著它們從一開始就必然履行相同的功能。為了更清楚地看到這一點，讓我們這樣思考。子宮的功能是懷孕；男性肢體的功能是使懷孕。但是，雖然這些肢體注定要履行這樣的功能，但並不因此它們從一開始就履行它們（因為我們看到許多女人不懷孕，如那些不生育的，即使她們有子宮），所以懷孕並非擁有子宮的直接和必然結果；但那些甚至不生育的人也避免性交，有些從一開始就是貞女，有些從某個時間開始。我們也看到男人保持童貞，有些從一開始，有些從某個時間開始；如此，藉著他們，因情慾而成為無法無天的婚姻被摧毀。我們發現即使是一些低等動物，雖然有子宮，也不生育，如騾子；公騾也不繁衍後代。因此，在人和非理性動物身上，我們都能看到性交被廢除；而這是在未來世界之前。我們的主耶穌基督由童貞女所生，沒有別的理由，只是為摧毀因無法無天的慾望而生育，並向統治者表明，天主無需人的干預就能創造人。當祂出生並接受了肉身的其他條件——我指的是飲食、衣著——只有這一履行性功能的條件祂沒有接受；因為，關於肉身的慾望，祂接受了一些必要的，而其他不必要的，祂沒有接受。因為如果肉身沒有飲食、衣著，它就會毀滅；但若沒有無法無天的慾望，它不會受到傷害。同時，祂預言在未來世界中，性交將被廢除；如祂所說：\BibleQuote{這世界的人有娶有嫁；但那來世的人既不娶也不嫁，卻要像天使一樣在天上。}\BibleRef{路 20:34} 不要讓那些不信的人驚奇，如果祂在來世廢除那些即使在今生也被廢除的肉體肢體的行為。\par

% structure: p0008 source=h2 target=section reason=relative-heading-rank; base=h2

\section{第四章 畸形的人必須以畸形復活嗎？}

嗯，他們說，如果肉身復活，它必須像它倒下時那樣復活；所以如果它死時只有一隻眼，它必須復活成一隻眼；如果瘸腿，就瘸腿；如果身體任何部分有缺陷，這個人必須在該部分有缺陷地復活。他們心靈的眼睛是多麼盲目！因為他們沒有在地上看見盲人重見光明，瘸子因祂的話行走。救主所做的一切，祂首先做是為應驗先知書中關於祂所說的，\BibleQuote{那時瞎子的眼睛要開啟，聾子的耳朵要開通，}\BibleRef{依 35:5} 等等；但也是為引起相信：在復活中，肉身將完整復活。因為如果祂在地上治癒了肉身的疾病，使身體完整，那麼在復活中祂更會如此，使肉身完美而完整地復活。這樣，他們那些可怕的難題將被醫治。\par

% structure: p0010 source=h2 target=section reason=relative-heading-rank; base=h2

\section{第五章 肉身的復活並非不可能}

但是再次，在那些主张肉身没有复活的人中，有些人断言这是不可能的；另一些人则认为，鉴于肉身是如此卑贱可鄙，天主复活它并不合适；还有一些人则说，肉身起初并没有领受恩许。那么，首先针对那些说天主使其复活是不可能的人，在我看来，我应当表明他们是无知的，他们口头上自称信徒，却以行为证明自己是不信者，甚至比不信者更不信。因为，既然所有外邦人都信他们的偶像，并相信对它们来说一切皆有可能（正如他们的诗人荷马甚至也说：\Quote{神明能行万事，且轻而易举；}他还加上了\Quote{轻而易举}这个词以凸显神明能力的伟大），许多人确实似乎比他们更不信。因为如果外邦人信他们的神，即偶像（\BibleQuote{有耳却听不见，有眼却看不见}），相信它们能行万事，尽管它们只是魔鬼，如圣经所说：\BibleQuote{列邦的神明是魔鬼}，那么，我们这些拥有正确、卓越、真实信德的人，更该相信我们的天主，因为我们也有（祂能力的）证明，首先在第一个人被创造时，因为他是由天主从泥土造成的；这足以证明天主的能力；接着，观察事物的人能看到人如何相互繁衍，并能更大程度地惊叹：从一小滴液体竟形成了如此伟大的活物。当然，如果这仅仅记载在恩许中而未见成就，那也会比另一件事更难以置信；但它因成就而变得更为可信。然而，就连复活的事，救主也向我们展示了成就，我们稍后将说到这些。不过现在我们正论证肉身复活的可能，请教会的子女原谅，如果我们援引看似世俗和自然的论据：第一，因为在天主面前，没有什么是世俗的，连世界本身也不是，因为它是祂的化工；第二，因为我们是在为应对不信者而进行论证。因为如果我们与信徒辩论，只说\Quote{我们相信}就够了；但现在我们必须以论证进行。上述证明确实足以表明肉身复活的可能性；但由于这些人极其不信，我们将进一步援引一个更有说服力的论据——一个并非取自信德的论据，因为他们不在信德范围内，而是取自他们自己的母亲——不信，当然，我指的是取自自然界理由。因为如果我们通过这样的论据向他们证明肉身复活是可能的，而他们既不能被信德的教导说服，也不能被世界的论据说服，那么他们当然值得极大的轻蔑。\par

% structure: p0012 source=h2 target=section reason=relative-heading-rank; base=h2

\section{第六章 复活与哲学家的意见相符}

那些被称为自然哲学家的人中，有些如\Person{柏拉图}，说宇宙由物质和天主构成；另一些如\Person{伊壁鸠鲁}，说由原子和虚空构成；还有些如斯多亚派，说由火、水、气、土这四元素构成。只需提及最流行的观点就够了。\Person{柏拉图}说，万物由天主按祂的设计从物质中造出；\Person{伊壁鸠鲁}及其追随者说，万物由原子和虚空通过某种身体自然运动的自行调节作用而造出；斯多亚派说，万物由四元素构成，天主贯穿其中。尽管他们之间存在如此分歧，但有些教义是他们共同承认的，其中之一就是：任何事物都不能从非存在中产生，任何事物也不能消灭或分解为全然非存在，并且元素是不灭的，万物由之生成。既然如此，按照所有这些哲学家的观点，肉身的再生将是可能的。因为如果按\Person{柏拉图}所说，是物质和天主，那么两者都是不灭的，且天主是工匠，即陶工的地位；物质则占据粘土或蜡或类似东西的地位。那么由物质形成的东西，无论是像还是雕像，是可灭的；但物质本身是不可灭的，如粘土、蜡或任何其他这类物质。这样，艺术家在粘土或蜡中设计，造出活物的形状；如果他的作品被毁，他并非不可能通过加工同样的材料并重新塑造而造出同样的形状。因此，按\Person{柏拉图}的说法，天主本身是不灭的，又有不灭的材料，即便最初造出的东西被毁，祂也可以重新造它，再造出与先前完全相同的形状。即使按斯多亚派的说法，身体由四种基本物质的混合产生，当这身体分解为四元素后，这些元素保持不灭，那么它们有可能第二次接受同样的融合和组合，由贯穿其中的天主来重新造出它们先前所形成的身体。好比一个人用金、银、铜、锡制成一种混合物，然后又想使其溶解，使各元素单独存在，再行混合，他若愿意，可以造出与先前完全相同的混合物。同样，按\Person{伊壁鸠鲁}的说法，原子和虚空是不灭的，通过原子聚合时的特定排列和调整，产生了一切其他构成物，包括身体本身；身体随时间分解，又分解为原先由之产生的那些原子。既然这些原子保持不灭，那么它们通过再次聚合、接受同样的排列和位置，造出与先前由它们产生的身体性质相同的身体，这绝非不可能。好比一个珠宝匠用马赛克拼出动物的形状，因时间或制作者本人的缘故，石头散落，他仍拥有散落的石头，可以重新收集，再以同样的方式排列，造出同样的动物形状。难道天主不能重新收集分解的肉身肢体，造出与先前由祂产生的相同的身体吗？\par

% structure: p0014 source=h2 target=section reason=relative-heading-rank; base=h2

\section{第七章　身体在天主眼中是珍贵的}

对于世俗之人，我已经充分证明了肉身复活的可能。如果连不信者的原则都不能发现肉身复活是不可能的，那么它更应合乎信徒的心意！按我们的顺序，现在必须对那些轻视肉身、说它不配复活和不配天上的安排的人说话，因为首先，它的本质是尘土；此外，因为它充满了各种邪恶，以致迫使灵魂与它一同犯罪。但这些人似乎不知道天主的一切工作，既不知道起初人的产生和形成，也不知道世上万物为何被造。因为圣言不是说了：\Quote{让我们按我们的肖像，照我们的模样造人？}\BibleRef{创 1:26}吗？什么样的人？显然祂指的是有肉身的人。因为圣言说：\Quote{天主用地上的尘土造了人}\BibleRef{创 2:7}。因此，显然按天主的肖像所造的人是有肉身的。那么，说由天主按祂自己的肖像所造的肉身是卑贱的、毫无价值，岂不荒谬吗？但肉身是天主珍贵的产业，这一点是明显的：首先，它由祂形成，如果至少肖像对造像者和艺术家是珍贵的；此外，它的价值可以从其他受造物的创造中看出。因为其他万物所为之而被造的那东西，对造物主来说是最珍贵的。\par

% structure: p0016 source=h2 target=section reason=relative-heading-rank; base=h2

\section{第八章　身体是否导致灵魂犯罪？}

他们说：很对，但肉身是罪人，以致迫使灵魂与它一同犯罪。这样他们就徒劳地控告它，并把两者的罪都归咎于它。但若无灵魂先行并激发它，肉身怎能独自犯罪？因为如同轭下的两头牛，若解开其中一头，两者都不能单独耕地；同样，若灵魂与身体解开了它们的结合，两者都不能单独成就任何事。如果肉身是罪人，那么救主唯独为它而来，如祂所说：\Quote{我不是来召义人，而是召罪人悔改}\BibleRef{谷 2:17}。既然肉身已被证明在天主眼中是珍贵的，在祂一切化工之上光荣的，那么它被祂拯救是极其正当的。\par

因此，我们必须面对那些说，即使肉体是天主特殊的化工，且为祂所珍视超越一切，但这并不立即意味着它拥有复活的应许。然而，难道这不荒谬吗？那以如此方式被造、且超越一切宝贵的，竟被其造主如此忽视，以致归于虚无？那么，雕刻家和画家，如果他们希望自己的作品存留，以从中获得荣耀，就会在作品开始朽坏时予以修复；但天主却会如此忽视祂自己的产业和化工，任其湮灭，不复存在。我们岂不应称此为徒劳？好比一个人建了房屋，却立即拆毁，或虽见其倾颓且能修复，却置之不理：我们会责备他徒劳无功；难道我们不应同样责备天主吗？但那位不朽者并非如此，宇宙的理智并非无感觉。让不信者缄默，即使他们自己不信。\par

但事实上，祂甚至已呼召肉体复活，并应许它永生。因为凡祂应许拯救人的地方，祂就把这应许给了肉体。因为人若不是由灵魂和肉体组成的理性动物，又是什么呢？灵魂本身是人吗？不，它是人的灵魂。肉体能被称为人吗？不，它被称为人的肉体。那么，如果这两者各自都不是人，而是两者结合而成的才被称为人，并且天主呼召人进入生命和复活，祂所呼召的就不是一部分，而是整体，即灵魂和肉体。因为，难道不是毫无疑问地荒谬吗？如果这两者在同一存在中、按照同一法则，一个得救而另一个不得救？而且，如果肉体再生并非不可能（正如已证明的），那么基于什么区别，灵魂得救而肉体不得救？他们使天主成为吝啬的天主吗？但祂是善的，并愿意所有人都得救。因着天主和祂的宣告，不仅你的灵魂（连同肉体）听到了并相信了耶稣基督，而且两者都受了洗，并行义。因此，他们使天主成为忘恩负义和不义的神，如果两者都信靠祂，祂却只愿拯救一个而不拯救另一个。好吧，他们说，但灵魂是不朽的，是天主的一部分，由祂所嘘，因此祂愿意拯救那属于自己的、与祂相似的部分；但肉体是可朽的，并非来自祂，如灵魂那样。那么，如果祂只立意拯救那在本性上得救的、作为祂一部分而存在的东西，这算是什么恩典？什么能力与善的彰显？因为灵魂从自身获得救恩；这样，拯救灵魂，天主并没有做什么大事。因为得救是它本性的命运，因为它是祂的一部分，是祂的嘘气。但拯救自己所有的，算不得什么恩典；因为这等于拯救自己。因为拯救自己一部分的人，是凭自己的手段拯救自己，以免那部分有所缺损；这不是一个善人的行为。因为即使一个人善待自己的子女后代，也不称他为善人；因为即使最凶猛的野兽也会这样做，而且为了幼崽，甚至甘愿赴死。但如果一个人为自己的奴隶做同样的事，那人才能公正地被称为善人。因此，救主也教导我们要爱仇敌，因为祂说：\Quote{你们有什么可夸的呢？}这样祂向我们表明，善举不仅在于爱那些从祂所生的人，也在于爱那些外人。而祂所吩咐我们的，祂自己首先做了。\par

% structure: p0020 source=h2 target=section reason=relative-heading-rank; base=h2

\section{第九章  基督的复活证明肉体复活}

如果祂不需要肉体，为什么祂医治它？而最有力的是，祂复活了死人。为什么？难道不是为了表明复活应该是怎样的吗？那么祂如何复活死人？他们的灵魂还是肉体？显然是两者。如果复活只是属灵的，那么祂在复活死人时，应该让肉体单独躺在一旁，而灵魂单独活着。但现在祂没有这样做，而是复活了肉体，在肉体中确认了生命的应许。为什么祂以受苦的肉体复活，除非为了显示肉体的复活？并且为了确认这一点，当祂的门徒不确定是否相信祂真的在肉体中复活了，看着祂而疑惑时，祂对他们说：\Quote{你们还没有信德，看看是我；}\BibleRef{路 24:32}祂让他们触摸祂，并向他们展示祂手上的钉痕。当他们通过种种证据确信是祂本人，且在肉体中时，他们请祂与他们一同进食，以便更准确地确知祂确实肉身复活了；祂吃了蜜和鱼。当祂这样向他们表明确实有肉体的复活时，为了也向他们表明这一点：肉体升天并非不可能（正如祂曾说过我们的居所在天上），\Quote{祂被接升天，他们亲眼看见，}\BibleRef{宗 1:9}祂是带着肉体升天的。因此，如果说了这一切之后，还有人要求关于复活的证明，他与撒杜塞人毫无区别，因为肉体的复活是天主的德能，超越一切推理，由信德所确立，并在工程中显现。\par

% structure: p0022 source=h2 target=section reason=relative-heading-rank; base=h2

\section{第十章  肉体得救，因此将复活}

复活是那死去之肉身的复活。因为精神不死；灵魂在肉身内，没有灵魂肉身就不能存活。肉身当灵魂离开它时，就不存在了。因为肉身是灵魂的房屋，而灵魂是精神的房屋。这三者，在一切对天主怀有真诚希望和无疑信心的人内，都将得救。因此，即使考察适合此世的种种论证，发现即使按照它们，肉身再生也并非不可能；而且，除了所有这些证据之外，救主在整个\WorkTitle{福音}中表明肉身有救恩，我们为何还要忍受那些不信的、危险的论证，为何看不出当我们听从这样的论点——说灵魂是不朽的，但肉身是会死的，不能复活——时，我们是在倒退呢？因为甚至在认识真理之前，我们就从\Person{毕达哥拉斯}和\Person{柏拉图}那里听到过这种说法。那么，如果救主说了这话，并只宣告灵魂得救，那么祂给我们带来了什么新东西，超越我们从\Person{毕达哥拉斯}、\Person{柏拉图}以及他们整个学派所听到的？但实际上，祂来是为向人宣布一个崭新而奇异的希望的喜讯。因为确实，天主应许祂不会在不朽中保持不朽，而要使朽坏变为不朽，这是一件奇异而新的事。但由于邪恶之君无法以其他方式败坏真理，就派出了他的宗徒（那些引入毒害教义的恶人），从那些把我们的救主钉在十字架上的人中拣选他们；这些人虽有救主的名号，却行那派遣他们者的作为，借着他，那名号本身也被毁谤了。但如果肉身不复活，为何还要守护它？为何我们不让它放纵欲望？为何我们不效法医生？据说，当医生遇到一个绝望且不治的病人时，就让他随意享受欲望。因为他们知道他快要死了；而那仇恨肉身的人确实这样做，尽可能地将肉身逐出它的产业；正因如此，他们也轻视它，因为它很快就要成为一具尸体。但如果我们的医生基督——天主，将我们从欲望中拯救出来后，以祂自己智慧而节制的法则来规整我们的肉身，那么显然祂保护它免受罪恶，因为它有得救的希望，正如医生不会让他们希望救活的病人随意享受快乐一样。\par

\SourceRecord{Translated by Marcus Dods.；Ante-Nicene Fathers；Vol. 1.；Edited by Alexander Roberts, James Donaldson, and A. Cleveland Coxe.；Buffalo, NY: Christian Literature Publishing Co.,；1885.；<http://www.newadvent.org/fathers/0131.htm>.}
